\documentclass[11pt]{article}

% ---------- Encoding & fonts ----------
\usepackage[T1]{fontenc}
\usepackage[utf8]{inputenc}
% Default Computer Modern to match the reference CV.
% For a Times-like look instead, uncomment: \usepackage{mathptmx}

% ---------- Geometry ----------
\usepackage[left=1.85in,right=0.8in,top=0.9in,bottom=1.05in]{geometry}

% ---------- Packages ----------
\usepackage{titlesec}
\usepackage{enumitem}
\usepackage{fancyhdr}
\usepackage[hidelinks]{hyperref}
\usepackage[spanish,es-noshorthands]{babel}
\usepackage{microtype}

% ---------- Editable metadata ----------
\newcommand{\cvname}{Santiago Rodríguez Vázquez}
\newcommand{\cvdate}{Julio de 2026}

% ---------- Marginal section headers (the Bugni look) ----------
\titleformat{\section}[leftmargin]
  {\normalfont\bfseries\filright}
  {}{0pt}{}
% {left width} {space above} {space below}
\titlespacing{\section}{1.35in}{1.3em}{0.6em}

% ---------- Footer: date left, page right ----------
\pagestyle{fancy}
\fancyhf{}
\renewcommand{\headrulewidth}{0pt}
%\fancyhf[R]{\footnotesize \cvdate}
\fancyfoot[C]{\footnotesize \thepage}

% ---------- Spacing helpers ----------
\setlength{\parindent}{0pt}
\newcommand{\entrydate}[1]{\hfill\textnormal{#1}}
\newlist{cvitems}{itemize}{1}
\setlist[cvitems]{leftmargin=1.2em,label=\textendash,itemsep=1.5pt,topsep=2pt,parsep=0pt}

% subgroup header inside a section (e.g. an institution)
\newcommand{\grp}[1]{\smallskip{\itshape #1}\par\nopagebreak\vspace{1pt}}

\begin{document}

% ================= HEADER =================
{\LARGE\bfseries {\cvname}}\par
\vspace{6pt}

\section{Contacto}
\begin{minipage}[t]{2.7in}
\href{mailto:santiago@rodriguez-vazquez.com.ar}{santiago@rodriguez-vazquez.com.ar}\\
\href{https://sanarovaz.github.io}{sanarovaz.github.io}
\\
+54 3442 516481
\end{minipage}


% ================= EDUCATION =================
\section{Educación}
\textbf{Universidad de San Andrés}.\entrydate{2026--Pres.}\\
Maestría en Economía.
\begin{cvitems}
  \item Núcleo aprobado. Optativas en curso: Comercio Internacional, Macroeconomía Dinámica,
        Métodos Econométricos para la Organización Industrial.
\end{cvitems}

\smallskip
\textbf{Universidad de Buenos Aires, Facultad de Cs. Económicas}.\entrydate{2019--2025}\\
Licenciatura en Economía.
\begin{cvitems}
  \item Promedio parcial: 8,03. Promedio real: 7,57.
  \item Materias del programa de Actuario: Estadística Actuarial, Análisis Numérico, Biometría.
  \item Tesina: \emph{``Determinantes del financiamiento de startups espaciales: proyectos de nanosatélites 2014--2024''} (nota: 10).
\end{cvitems}

% ================= RESEARCH EXPERIENCE =================
\section{Experiencia en investigación}
\textbf{Instituto Interdisciplinario de Economía Política (IIEP)}.\entrydate{2025}\\
Becario UBACyT Estímulo.
\begin{cvitems}
    \item  Proyecto \emph{``Narcoeconomía en Argentina: dimensionamiento, impactos y propuestas de política''}, a cargo de Andrés López (UBA).
\end{cvitems}

% ================= TEACHING =================
\section{Experiencia docente}
\textbf{Universidad de San Andrés}.
\begin{cvitems}
  \item Tutor, Economía Internacional (E155).\entrydate{2026--Pres.}
\end{cvitems}

\textbf{Universidad de Buenos Aires, Facultad de Cs. Económicas}.
\begin{cvitems}
  \item Ayudante de Primera, Microeconomía I (291).\entrydate{2026--Pres.}
  \item Docente Adscripto, Desarrollo Económico (559). \entrydate{2025--Pres.}
  \item Ayudante de Primera, Economía Financiera (549).\entrydate{2025--Pres.}
  \item Ayudante de Primera, Microeconomía II (286).\entrydate{2025--Pres.}
  \item Ayudante de Segunda, Microeconomía II (286).\entrydate{2022--2025.}
  \end{cvitems}

% ================= PROFESSIONAL EXPERIENCE =================
\section{Experiencia profesional}
\textbf{Veritran}.\entrydate{2022--2023}\\
Finance Analyst Jr.
\begin{cvitems}
  \item Apoyo en cuentas a pagar para subsidiarias de Latinoamérica, EE.UU. y Europa.
  \item Mantenimiento de bases de datos de empleados y proveedores en NetSuite.
  \item Contacto con proveedores, seguimiento de procesos de actualización tributaria e
        implementación de sistemas internos para órdenes de compra.
\end{cvitems}

% ================= SKILLS =================
\section{Capacidades}
\textbf{Técnicas}
\begin{cvitems}
  \item Programación: R, Python. Familiar con SQL, Stata, Tableau Public.
  \item Paquetes y presentaciones: LaTeX, Google Workspace, Microsoft Office.
\end{cvitems}
\textbf{Idiomas}
\begin{cvitems}
\item Inglés (avanzado): Cambridge \emph{Certificate in Advanced English}, puntaje 198, 2017.
\end{cvitems}

% ================= SEMINARS =================
\section{Participación en eventos CyT}

\textbf{Ciclo de seminarios permanentes del Departamento de Economía, Universidad de San Andrés}
\begin{cvitems}
  \item Bugni, F. (2026). \emph{Identification and Inference on Treatment Effects under Covariate-Adaptive Randomization and Imperfect Compliance}.
  \item Caravello, T. (2026). \emph{Pricing Mistakes and the Non-Neutrality of Money}.
  \item Tobal, M. (2026). \emph{Monetary Policy Transmission in an Emerging Market: The Financial- Friction Channel vs. The Interest-Rate Channel}.
  \item Ballestero, G. (2026). \emph{Centralized Admission Mechanisms and International Migration}.
  \item Amarante, V. (2026). \emph{Reproducing Gender Inequality: Adolescent Time Use and Parental Influence in Colombia and Mexico}.
  \item Giulla, T. (2026). \emph{What Candidates Do Parties Pick? Executive Experience and Candidate Selection}.
  \item Satyanath, S. (2026). \emph{Social Clubs and Political Office}.
  \item Roldán, F. (2026). \emph{The Perils of Bilateral Soverign Debt}.
\end{cvitems}

\vspace{5px}
\textbf{Seminarios del Instituto Interdisciplinario de Economía Política (IIEP)}
\begin{cvitems}
  \item Edo, M. (2025). \emph{Gender Gaps in Academic Careers: Evidence from Economics in Argentina.}
  \item Maurizio, R. (2025). \emph{Short- and long-term earning effects of job exits: insights from matched employer-employee data in Argentina.}
  \item Aromí, D. (2025). \emph{Technological change and labor outcomes: measuring the role of automation and augmentation.}
  \item Trombetta, M. (2025). \emph{Intergenerational occupational transmission as a determinant of social mobility.}
  \item Escolar, J.\,M. (2025). \emph{Bidding for Firms: Efficiency vs.\ Redistribution and Optimal Government Policy.}
  \item Bisang, R. (2024). \emph{Aportes de la bioeconomía al desarrollo argentino: ¿nueva oportunidad de la ``industrialización biológica''?}
  \item Viollaz, M. (2024). \emph{Women Political Leaders as Agents of Environmental Change.}
  \item Petralia, S. (2024). \emph{Innovation Complementarities in Open Source Software.}
  \item Tomassi, M. (2024). \emph{The political economy of redistribution and (in)efficiency in Latin America and the Caribbean.}
  \item Díaz de Astarloa, B. (2024). \emph{Evolución reciente del dinamismo empresarial y la productividad en la industria manufacturera argentina.}
  \item Depetris Chauvin, N. (2024). \emph{Promoting Exports in Brazil.}
  \item You, H.\,Y. (2024). \emph{How Are Politicians Informed? Witness and Information Provision in Congress.}
\end{cvitems}

\vspace{5px}
\textbf{Seminario Interuniversitario sobre Desarrollo Productivo Argentino (SIDPA)}
\begin{cvitems}
  \item López, A., Sturbin, L.\ \& Pascuini, P. (2025). \emph{Saliendo al mundo. Internacionalización de startups IA y Biotecnología en Argentina.}
  \item Bentivegna, B. (2025). \emph{Análisis del estado actual de los sectores de soja y carne vacuna en Argentina a la luz del Reglamento (UE) 2023/1115 sobre productos libres de deforestación.}
  \item García Díaz, F. (2024). \emph{Asimetrías y desigualdades territoriales en la Argentina. Aportes para el debate.}
\end{cvitems}

\end{document}
